\section{Related Work}

\noindent\textbf{Vision-based tactile sensing.}
GelSight-type sensors convert elastomer deformation into images under structured illumination~\cite{johnson2009retrographic,yuan2017gelsight}; GelSlim~\cite{donlon2018gelslim,sipos2024gelslim4} and DIGIT~\cite{lambeta2020digit} adapt the principle to fingertips. Because the measurement depends jointly on gel mechanics, LED layout and camera response, models transfer poorly between sensors; sensor-invariant representations~\cite{gupta2025sitr,zhao2024transferable,higuera2024sparsh,higuera2025sparshx} and deformation-space translation~\cite{zai2025across,yang2025flow} attack this on the representation side. We use a frozen sensor-invariant encoder~\cite{gupta2025sitr} as backbone so that our comparison isolates the data.

\noindent\textbf{Tactile simulation.}
Physically based tactile renderers synthesize GelSight images from mesh contacts~\cite{agarwal2021simulation,wang2022tacto,luu2023simulation}, with example-based~\cite{si2022taxim}, particle-based~\cite{chen2023tacchi}, differentiable~\cite{xu2022efficient} and GPU-scale~\cite{nguyen2024tacex,akinola2025tacsl} variants. Learned renderers instead map simulated depth or contact state to sensor images with an image-to-image network~\cite{isola2017pix2pix,church2022tactile,zhong2025tactgen}. Our two renderers follow the two lines and share meshes, contact sampler and labels, so they differ only in the appearance model.

\noindent\textbf{Sim-to-real transfer and its evaluation.}
Domain randomization~\cite{tobin2017domain,peng2018sim} and structured randomization around a calibrated center~\cite{chen2024general} trade fidelity for robustness; sensor-background conditioning~\cite{zhong2025tactgen} and low-fidelity pretraining~\cite{gano2025lowfidelity} report gains from synthetic tactile data. These works measure the benefit of simulation but rarely report how much real data the simulator consumed, or whether the simulator saw the test distribution. Cross-sensor transfer is usually attacked on the representation side~\cite{gupta2025sitr,zhao2024transferable}; we instead ask how much of a new sensor's or session's real data a renderer needs. Our contribution is the evaluation protocol---renderer and downstream model trained on the same disjoint real subset, tested on held-out objects, a second session and a second sensor---and the resulting data-efficiency curves for depth, normals and pose with both renderer families.
